%% notes-tex/database-expansion-100/sections/rule.tex
%% [[experiments.database.expansion-100]]
%% https://github.com/Mjvolk3/torchcell/tree/main/notes-tex/database-expansion-100/sections/rule.tex
%% SOURCE: hand-authored prose; every number quoted here comes from
%% experiments/database/scripts/build_candidate_datasets_table.py

\section{The selection rule\secstatus{todo}}\label{sec:rule}

\subsection{Sequence reconstruction is the gate, not a preference\secstatus{todo}}

A phenotype is only trainable against a genotype that can be written down. Every row in
Table~\ref{tab:candidates} therefore names a \term{sequence basis}: the route by which the
total genomic content of one strain would be reconstructed. The eleven bases in use are

\begin{itemize}[leftmargin=1.2em,itemsep=1pt,topsep=2pt]
  \item \term{S288C-KO} -- the reference genome minus one cataloged ORF (BY4741, BY4742,
    BY4743). Thirty of the 150 that reach 200 sit here, the largest single group.
  \item \term{S288C-KO/het} -- a heterozygous diploid deletion, so a dosage edit rather
    than a null.
  \item \term{S288C+guide} -- the genome is unedited and the perturbation is a designed
    cassette plus a guide target. Regulatory, not sequence-level.
  \item \term{S288C+designed-edit} -- a designed single-nucleotide or small-indel library.
    Designed, and not verified per strain.
  \item \term{S288C+tag} -- the reference plus a designed tag, degron or promoter cassette.
  \item \term{S288C+ORF-plasmid} -- the reference plus a barcoded ORF on a known plasmid.
  \item \term{S288C+reporter-locus} -- the reference plus a designed regulatory sequence
    driving a reporter, which is what every massively parallel reporter assay perturbs.
  \item \term{isolate-WGS} -- a published assembly or variant call set per isolate.
  \item \term{segregant-WGS} -- genotype calls or whole-genome sequence per progeny.
  \item \term{engineered-chassis} -- a named production strain plus heterologous cassettes.
  \item \term{reference-only} -- wild type, where the environment carries the perturbation.
\end{itemize}

What fails the gate is a strain whose genome is unknown even in principle. Genome shuffling
and adaptive laboratory evolution without resequencing both land there: the selected
progeny differ from their parents at unmapped positions, so there is no genotype to map a
phenotype from. Those are dropped in Table~\ref{tab:excluded} rather than ranked low, since
no amount of phenotypic value substitutes for a missing genotype.

The \term{S288C+designed-edit} and \term{S288C+guide} bases carry a weaker claim than the
others and the distinction is worth holding. A designed variant library states what each
member was built to be, not what it turned out to be, so its genotypes are designed rather
than realized. Ingesting one commits to storing the design and a verification status, not a
verified sequence.

\subsection{Band comes before tier, and why\secstatus{todo}}\label{sec:band}

Scale is the right default ordering and it is the wrong first question when a specific
campaign is being planned. The largest row in this table, de Boer 2020 at $10^8$ promoter
sequences, tells a Perturb-seq design nothing; a 61,094-guide CRISPR interference library
with a released per-guide matrix tells it what guides to order. Each row therefore carries
a \term{band}, applied before tier and before scale:

\begin{itemize}[leftmargin=1.2em,itemsep=1pt,topsep=2pt]
  \item \term{perturb-seq}, 30 rows. A CRISPR interference or activation library whose
    per-guide matrix is released, a barcoded genotype set read out per cell, a designed
    single-gene perturbation crossed with a transcriptome, or the strain and noise
    measurements a pooled per-cell campaign depends on.
  \item \term{molecular layers}, 23 rows. Two or more of transcript, translation,
    protein, phosphosite, metabolite, flux and turnover, joinable on one axis, either
    inside one study or against a supported dataset. The join is per \emph{strain} when
    both layers are measured on the same genotypes, and per \emph{gene} when one layer is
    a single genome-wide coefficient such as a half-life.
  \item \term{scale}, 115 rows. Everything else, ordered as before.
\end{itemize}

The second band was \term{metabolism x expression} and is widened, because transcript
against protein is the same join as expression against flux, and because the rows that say
\emph{why} transcript and protein disagree, translation efficiency and turnover, had no
band to sit in at all. Its membership is unchanged apart from the rows added with it.

Within a band the order is tier, then measurements. The cost is visible and is not hidden:
the two bands now total 53 rows, so they fill wave 1 and three of them spill into wave 2,
and no scale row reaches wave 1 at all. de Boer 2020 falls from 2 to 54 and Lee 2014 from 4
to 56. Both are the largest things in the table by instance count and both are now second
wave, which is the deliberate consequence of ordering by what a row combines with rather
than by how many numbers it holds. A reader who wants scale back reads wave 2 first. Every
move is listed in Table~\ref{tab:swaps}.

\subsection{Tier, and what rows are ranked on\secstatus{todo}}

Tier is assigned from stated bars rather than judgment:

\begin{itemize}[leftmargin=1.2em,itemsep=1pt,topsep=2pt]
  \item \textbf{Tier 1} clears all three of $\geq$$10^3$ genotypes, $\geq$$10^4$ instances,
    and a clean sequence basis, and lands in one of the nine classes.
  \item \textbf{Tier 2} clears two of the three, or is the only dataset covering a
    requested class, or directly de-risks the Perturb-seq proposal.
  \item \textbf{Tier 3} is a modality or reference backbone the substrate needs without
    itself carrying a large genotype axis.
  \item \textbf{Tier 4} is real but small, coarsely labeled, or expensive to extract.
\end{itemize}

The bars apply to what a row contains rather than to its reputation. Piotrowski 2017 has
the largest compound library in yeast at 13,524 compounds and still fails Tier 1, because
157 strains is not a genotype axis. Turco 2023 has the widest environment axis at 7,536 and
fails for a different reason: it is a meta-aggregation, so its instance count re-serves
primary screens, several of them already built, and the net-new fraction is unknown until
it is de-duplicated against the rest of this table.

Within a tier, rows sort by \term{measurements}, meaning instances times phenotype
dimensionality. Instances alone is the wrong statistic and gets the order visibly wrong:
Muenzner 2024 is 796 instances, which ranks it below a mid-sized fitness screen, but each
instance is a proteome, so it is roughly $1.6\times10^6$ numbers. This is the same
normalization the supported-datasets table already applies through its gzip-signal column.

\subsection{Where the lines fall, and what they cost\secstatus{todo}}

The long-run cut at 150 is arithmetic, $200-50$, not a quality threshold. Rows on both
sides of it are close, so it should be read as a budget rather than a verdict.

The two nearer lines are what a build week is planned against. Wave 1 is rows 1--50 and
wave 2 rows 51--70; a wave-2 row is promoted the moment a wave-1 row turns out to have no
recoverable per-strain data, which has already happened once to Lee 2014. Wave 1 is 30
Perturb-seq rows and 20 molecular-layer rows, with no scale row in it, and wave 2 is three
molecular-layer rows followed by scale, so the two waves differ in kind and not only in
rank.

A dataset named explicitly in the scoping request is pinned into the recommended set even
if it ranks below the cut, and the pin is reported rather than folded into the score, so
the set stays auditable. No pin is currently in force: Jackson 2020, the one requested row,
ranks 3 unaided.

Displacement, when a pin does bind, picks the weakest non-requested row by tier first and
measurements second. The tier term is load-bearing rather than decorative: ranking on
measurements alone selected Puddu 2019 for eviction, and that row is the whole-collection
sequencing on which every \term{S288C-KO} sequence basis in this table rests. Banding has
since moved Puddu to 4 on its own merits, because a pooled campaign reads a barcode and
infers a genotype and this is the measurement of how often that inference is wrong.
